\chapter{第1章\texorpdfstring{\hspace{0.5em}}{ }质点运动学}

质点运动学（particle kinematics）的任务是描述作机械运动的物体在空间的位置随时间变化的关系，而不涉及引起运动和改变运动的原因。本章先阐明描述物质运动的测量标准与方法，再定义位置矢量、位移、速度、加速度等描述质点运动的物理量，讨论它们随时间变化的规律；然后讨论曲线运动中的法向加速度和切向加速度，以及圆周运动的角量描述；最后介绍相对运动与速度叠加原理。

\section{运动\texorpdfstring{\hspace{0.5em}}{ }时空\texorpdfstring{\hspace{0.5em}}{ }测量}
    \subsection{物质与运动}
    \begin{enumerate}
        \item 物理学（physics）是研究物质的基本结构、基本相互作用和基本运动规律的科学。
        \item 物质以\textbf{实物}（material object）和\textbf{场}（field）两种基本形式存在：物质是客观存在，是能量的表达；运动是物质的固有属性，由能量的涨落引起。
        \item 真空是一切场的基态，实物与场都是真空的激发态。二者的主要区别在于：实物的能量密度远大于场的能量密度。场的激发状态表现为出现相应的粒子，粒子以一定方式聚焦起来就构成实物。在物质存在的两种基本形式中，场是更基本的。
        \item 质点运动学只回答“怎样动”，不回答“为什么动”（后者属动力学）。
    \end{enumerate}

    \subsection{测量的标准}
    \begin{enumerate}
        \item 物理学以测量为基础。物理量的\textbf{测量}，就是把待测物理量与一个公认的标准量相比较，所得比值即为数值，并配以约定的单位。
        \item 标准选取原则：\textbf{易于复现}、\textbf{恒定不变}（足够精确、稳定，与外界条件无关）。
    \end{enumerate}

    \subsection{国际单位制与量纲}
        \subsubsection{SI 的七个基本量}
        由基本量按定义或物理定律导出的物理量称为\textbf{导出量}，其单位为导出单位。
        \begin{center}
        \begin{tabular}{llll}
        \toprule
        基本量 & 量纲符号 & 单位名称 & 单位符号 \\
        \midrule
        长度 & $\mathrm{L}$ & 米 & $\mathrm{m}$ \\
        质量 & $\mathrm{M}$ & 千克 & $\mathrm{kg}$ \\
        时间 & $\mathrm{T}$ & 秒 & $\mathrm{s}$ \\
        电流 & $\mathrm{I}$ & 安培 & $\mathrm{A}$ \\
        热力学温度 & $\varTheta$ & 开尔文 & $\mathrm{K}$ \\
        物质的量 & $\mathrm{N}$ & 摩尔 & $\mathrm{mol}$ \\
        发光强度 & $\mathrm{J}$ & 坎德拉 & $\mathrm{cd}$ \\
        \bottomrule
        \end{tabular}
        \end{center}

        \subsubsection{量纲与量纲一致性原理}
        表示一个物理量由哪些基本量组合而成的式子称为该物理量的\textbf{量纲}（dimension）：
        \keyformula{\mathrm{dim}\,Q=\mathrm{L}^{\alpha}\mathrm{M}^{\beta}\mathrm{T}^{\gamma}\mathrm{I}^{\delta}\varTheta^{\varepsilon}\mathrm{N}^{\zeta}\mathrm{J}^{\eta}}

        \textbf{量纲一致性原理}：物理方程等式两侧的量纲必须完全相同；指数、对数、三角函数的变量必须是无量纲量。它的三个用途是：校验公式正误、单位换算、由量纲分析猜测物理量之间的函数关系。量纲指数全为零的量为\textbf{无量纲量}。

        \begin{center}
        \begin{tabular}{lll}
        \toprule
        物理量 & 单位 & 量纲 \\
        \midrule
        速度 & $\mathrm{m\cdot s^{-1}}$ & $\mathrm{LT^{-1}}$ \\
        加速度 & $\mathrm{m\cdot s^{-2}}$ & $\mathrm{LT^{-2}}$ \\
        力 & $\mathrm{N}$ & $\mathrm{MLT^{-2}}$ \\
        功、能量 & $\mathrm{J}$ & $\mathrm{ML^{2}T^{-2}}$ \\
        动量 & $\mathrm{kg\cdot m\cdot s^{-1}}$ & $\mathrm{MLT^{-1}}$ \\
        角动量 & $\mathrm{kg\cdot m^{2}\cdot s^{-1}}$ & $\mathrm{ML^{2}T^{-1}}$ \\
        转动惯量 & $\mathrm{kg\cdot m^{2}}$ & $\mathrm{ML^{2}}$ \\
        \bottomrule
        \end{tabular}
        \end{center}

    \subsection{时间与空间}
    \begin{enumerate}
        \item 时空是运动物质的存在形式，客观存在，不能脱离物质而独立存在；时间与空间通过物质运动相互关联。
        \item \textbf{时间}描述物质运动的持续性与顺序性。1967 年规定：铯-133 原子基态两个超精细能级跃迁辐射 $9\,192\,631\,770$ 个周期的持续时间为 $1\ \mathrm{s}$。
        \item \textbf{长度}描述物质的广延性。1983 年规定：光在真空中 $1/299\,792\,458\ \mathrm{s}$ 内传播的路径长度为 $1\ \mathrm{m}$。
        \item 经典力学默认“绝对时间、绝对长度”：时间的测量与长度（空间间隔）的测量均与参考系无关。这一结论只在低速（$v\ll c$）情形下近似成立，高速情形须由狭义相对论修正。
    \end{enumerate}

    \subsection{参考系和坐标系}
        \subsubsection{参考系}
        运动具有绝对性，而对运动的描述具有相对性。被选作观测基准的物体（或物体群）称为\textbf{参考系}（reference frame）。同一运动在不同参考系中的轨迹、位移、速度一般不同。
        \corepoint{同一问题的所有物理量（位矢、速度、加速度、力、能量）必须统一到同一参考系中计算，这是解运动学、动力学题的第一条纪律。}

        \subsubsection{坐标系}
        参考系选定后，还要固结一组有刻度的射线、曲线或角度以定量描述位置，这就是\textbf{坐标系}。常用：直角坐标系、自然坐标系、极坐标系、柱坐标系、球坐标系。坐标系可任选，选取原则是使运算最简。

\section{质点运动的描述}
    \subsection{质点}
    当物体的大小、形状对所研究的问题影响不大，或物体上各部分运动规律相同时，可把物体抽象为集中全部质量的几何点，即\textbf{质点}（particle）。质点是力学中的理想化模型，其“可否适用”是相对的：研究地球公转时可把地球视为质点，研究地球自转时则不能。

    \subsection{位置矢量与运动方程}
    由参考系原点指向质点所在位置的矢量称为\textbf{位置矢量}（简称位矢），记作 $\boldsymbol{r}$。直角坐标系中
    \keyformula{\boldsymbol{r}(t)=x(t)\boldsymbol{i}+y(t)\boldsymbol{j}+z(t)\boldsymbol{k}}
    位矢的大小与方向余弦为
    \[
    r=|\boldsymbol{r}|=\sqrt{x^{2}+y^{2}+z^{2}},\qquad
    \cos\alpha=\frac{x}{r},\ \cos\beta=\frac{y}{r},\ \cos\gamma=\frac{z}{r}
    \]
    \begin{enumerate}
        \item \textbf{运动方程}：$\boldsymbol{r}=\boldsymbol{r}(t)$，或其分量式 $x=x(t),\ y=y(t),\ z=z(t)$，描述位矢随时间变化的规律。
        \item \textbf{轨道方程}：从运动方程中消去参数 $t$，得到坐标间的函数关系 $f(x,y,z)=0$，即质点运动的轨迹。
    \end{enumerate}

    \subsection{位移与路程}
    \begin{enumerate}
        \item \textbf{位移} $\Delta\boldsymbol{r}$：$\Delta t$ 时间内位矢的增量，是矢量，方向由初位置指向末位置。
        \[
        \Delta\boldsymbol{r}=\boldsymbol{r}_{B}-\boldsymbol{r}_{A}
        =\Delta x\,\boldsymbol{i}+\Delta y\,\boldsymbol{j}+\Delta z\,\boldsymbol{k},
        \qquad |\Delta\boldsymbol{r}|=\sqrt{\Delta x^{2}+\Delta y^{2}+\Delta z^{2}}
        \]
        \item \textbf{路程} $\Delta s$：质点实际走过的轨迹长度，是标量，与路径有关。
        \item 一般情况下 $|\Delta\boldsymbol{r}|\neq\Delta s$，且 $|\Delta\boldsymbol{r}|\leqslant\Delta s$；只有单向直线运动或 $\Delta t\to0$ 时才有 $\mathrm{d}s=|\mathrm{d}\boldsymbol{r}|$。
    \end{enumerate}

    \subsection{速度矢量}
        \subsubsection{平均速度与平均速率}
        \[
        \bar{\boldsymbol{v}}=\frac{\Delta\boldsymbol{r}}{\Delta t},\qquad
        \bar{v}=\frac{\Delta s}{\Delta t}
        \]
        平均速度是矢量，平均速率是标量。\corepoint{平均速度的大小 $|\bar{\boldsymbol{v}}|$ 一般不等于平均速率 $\bar{v}$（绕行一周时位移为零而路程不为零）。}

        \subsubsection{瞬时速度与速率}
        瞬时速度（简称速度）是平均速度在 $\Delta t\to0$ 时的极限，方向沿轨迹在该点的切线并指向质点前进的一侧：
        \keyformula{\boldsymbol{v}=\lim_{\Delta t\to0}\frac{\Delta\boldsymbol{r}}{\Delta t}=\frac{\mathrm{d}\boldsymbol{r}}{\mathrm{d}t}
        =v_{x}\boldsymbol{i}+v_{y}\boldsymbol{j}+v_{z}\boldsymbol{k}}
        速率是速度的大小：
        \[
        v=|\boldsymbol{v}|=\frac{\mathrm{d}s}{\mathrm{d}t}=\sqrt{v_{x}^{2}+v_{y}^{2}+v_{z}^{2}},
        \qquad v_{x}=\frac{\mathrm{d}x}{\mathrm{d}t},\ v_{y}=\frac{\mathrm{d}y}{\mathrm{d}t},\ v_{z}=\frac{\mathrm{d}z}{\mathrm{d}t}
        \]

\section{质点运动变化的描述}
    \subsection{加速度矢量}
    加速度是描述速度矢量（大小与方向）随时间变化快慢的物理量。
    \keyformula{\boldsymbol{a}=\lim_{\Delta t\to0}\frac{\Delta\boldsymbol{v}}{\Delta t}
    =\frac{\mathrm{d}\boldsymbol{v}}{\mathrm{d}t}=\frac{\mathrm{d}^{2}\boldsymbol{r}}{\mathrm{d}t^{2}}
    =a_{x}\boldsymbol{i}+a_{y}\boldsymbol{j}+a_{z}\boldsymbol{k}}
    其中 $a=\sqrt{a_{x}^{2}+a_{y}^{2}+a_{z}^{2}}$，单位 $\mathrm{m\cdot s^{-2}}$。

    \subsubsection{质点运动学的两类问题}
    \begin{enumerate}
        \item \textbf{微分问题}：已知运动方程 $\boldsymbol{r}=\boldsymbol{r}(t)$，求导得速度、加速度。
        \item \textbf{积分问题}：已知 $\boldsymbol{a}(t)$（或 $\boldsymbol{a}(\boldsymbol{v})$、$\boldsymbol{a}(\boldsymbol{r})$）及初始条件，积分求 $\boldsymbol{v}(t)$、$\boldsymbol{r}(t)$。
        \[
        \boldsymbol{v}=\boldsymbol{v}_{0}+\int_{0}^{t}\boldsymbol{a}\,\mathrm{d}t,\qquad
        \boldsymbol{r}=\boldsymbol{r}_{0}+\int_{0}^{t}\boldsymbol{v}\,\mathrm{d}t
        \]
        当 $a=a(v)$ 时用 $\displaystyle a=\frac{\mathrm{d}v}{\mathrm{d}t}$ 分离变量；当 $a=a(x)$ 时用 $\displaystyle a=\frac{\mathrm{d}v}{\mathrm{d}t}=v\frac{\mathrm{d}v}{\mathrm{d}x}$ 换元积分。
    \end{enumerate}

    \subsubsection{运动叠加原理}
    任意曲线运动都可分解为三个相互正交方向上独立直线运动的叠加，各分运动互不影响。这是把复杂运动化为若干简单运动的依据（如平抛运动 = 水平匀速直线 + 竖直自由落体）。

    \subsection{自然坐标系中的速度和加速度}
    以轨迹上的切向单位矢量 $\boldsymbol{e}_{\mathrm{t}}$（指向质点前进方向）和法向单位矢量 $\boldsymbol{e}_{\mathrm{n}}$（指向曲率中心）构成的坐标系称为自然坐标系，适用于轨迹已知的平面曲线运动。

    \subsubsection{自然坐标下的速度}
    \[
    \boldsymbol{v}=\frac{\mathrm{d}s}{\mathrm{d}t}\boldsymbol{e}_{\mathrm{t}}=v\boldsymbol{e}_{\mathrm{t}}
    \]

    \subsubsection{切向加速度与法向加速度}
    对 $\boldsymbol{v}=v\boldsymbol{e}_{\mathrm{t}}$ 求导，注意 $\boldsymbol{e}_{\mathrm{t}}$ 的方向随位置改变，得
    \keyformula{\boldsymbol{a}=\frac{\mathrm{d}v}{\mathrm{d}t}\boldsymbol{e}_{\mathrm{t}}+\frac{v^{2}}{\rho}\boldsymbol{e}_{\mathrm{n}}
    =a_{\mathrm{t}}\boldsymbol{e}_{\mathrm{t}}+a_{\mathrm{n}}\boldsymbol{e}_{\mathrm{n}}}
    \begin{enumerate}
        \item \textbf{切向加速度} $a_{\mathrm{t}}=\dfrac{\mathrm{d}v}{\mathrm{d}t}$，反映速度\textbf{大小}的变化；
        \item \textbf{法向加速度（向心加速度）} $a_{\mathrm{n}}=\dfrac{v^{2}}{\rho}$，反映速度\textbf{方向}的变化，$\rho$ 为该点轨迹的曲率半径；
        \item 合加速度大小 $a=\sqrt{a_{\mathrm{t}}^{2}+a_{\mathrm{n}}^{2}}$，方向与速度的夹角 $\theta=\arctan\left(a_{\mathrm{n}}/a_{\mathrm{t}}\right)$。
    \end{enumerate}

    \begin{center}
    \begin{tabular}{ccc}
    \toprule
    切向分量 & 法向分量 & 运动类型 \\
    \midrule
    $a_{\mathrm{t}}=0$ & $a_{\mathrm{n}}=0$ & 匀速直线运动 \\
    $a_{\mathrm{t}}\neq0$ & $a_{\mathrm{n}}=0$ & 变速直线运动 \\
    $a_{\mathrm{t}}=0$ & $a_{\mathrm{n}}\neq0$ & 匀速率曲线运动（含匀速圆周运动） \\
    $a_{\mathrm{t}}\neq0$ & $a_{\mathrm{n}}\neq0$ & 变速率曲线运动 \\
    \bottomrule
    \end{tabular}
    \end{center}

    \subsection{圆周运动的角量描述}
        \subsubsection{角量与线量的定义}
        \begin{enumerate}
            \item \textbf{线量}：弧长 $s$、速率 $v$、切向加速度 $a_{\mathrm{t}}$、法向加速度 $a_{\mathrm{n}}$；
            \item \textbf{角量}：角位置 $\theta$、角位移 $\Delta\theta$、角速度 $\omega$、角加速度 $\alpha$。
        \end{enumerate}
        规定：逆时针转动取正，顺时针转动取负。

        \subsubsection{角速度与角加速度}
        \keyformula{\omega=\lim_{\Delta t\to0}\frac{\Delta\theta}{\Delta t}=\frac{\mathrm{d}\theta}{\mathrm{d}t},
        \qquad
        \alpha=\lim_{\Delta t\to0}\frac{\Delta\omega}{\Delta t}=\frac{\mathrm{d}\omega}{\mathrm{d}t}
        =\frac{\mathrm{d}^{2}\theta}{\mathrm{d}t^{2}}}
        单位分别为 $\mathrm{rad\cdot s^{-1}}$、$\mathrm{rad\cdot s^{-2}}$。

        \subsubsection{角量与线量的关系}
        \keyformula{v=R\omega,\qquad a_{\mathrm{t}}=R\alpha,\qquad a_{\mathrm{n}}=R\omega^{2}=\frac{v^{2}}{R}}

        \subsubsection{角速度与角加速度的矢量表达}
        无限小角位移满足矢量运算法则，角速度 $\boldsymbol{\omega}$、角加速度 $\boldsymbol{\alpha}$ 的方向沿转轴，由右手螺旋定则判定。
        \keyformula{\boldsymbol{v}=\boldsymbol{\omega}\times\boldsymbol{r},\qquad
        \boldsymbol{a}=\boldsymbol{\alpha}\times\boldsymbol{r}+\boldsymbol{\omega}\times\boldsymbol{v}}
        式中第一项为切向加速度，第二项为法向加速度。

    \subsection{典型运动的运动学模型}
        \subsubsection{匀变速直线运动}
        \keyformula{v=v_{0}+at,\qquad x=x_{0}+v_{0}t+\frac{1}{2}at^{2},\qquad v^{2}-v_{0}^{2}=2a(x-x_{0})}
        自由落体（$a=g$，初速为零）是 $x_{0}=0,\ v_{0}=0$ 的特例。

        \subsubsection{抛体运动}
        以初速率 $v_{0}$、仰角 $\theta$ 抛出，取抛出点为原点、$y$ 轴竖直向上，加速度 $a_{x}=0,\ a_{y}=-g$：
        \keyformula{x=v_{0}\cos\theta\cdot t,\qquad y=v_{0}\sin\theta\cdot t-\frac{1}{2}gt^{2}}
        消去 $t$ 得轨道方程（抛物线）
        \[
        y=x\tan\theta-\frac{gx^{2}}{2v_{0}^{2}\cos^{2}\theta}
        \]
        三个常用结论：
        \[
        T=\frac{2v_{0}\sin\theta}{g},\qquad
        H=\frac{v_{0}^{2}\sin^{2}\theta}{2g},\qquad
        R=\frac{v_{0}^{2}\sin2\theta}{g}
        \]
        其中 $T$ 为飞行时间，$H$ 为最大高度，$R$ 为射程。$\theta=45^{\circ}$ 时射程最大；$\theta$ 与 $90^{\circ}-\theta$ 的射程相同。平抛运动是 $\theta=0$ 的特例，斜抛的“运动叠加”分解为水平匀速直线运动与竖直匀变速直线运动。

        \subsubsection{匀变速圆周运动}
        以角量表示的公式与匀变速直线运动完全对应：
        \keyformula{\omega=\omega_{0}+\alpha t,\qquad
        \theta=\theta_{0}+\omega_{0}t+\frac{1}{2}\alpha t^{2},\qquad
        \omega^{2}-\omega_{0}^{2}=2\alpha(\theta-\theta_{0})}

        \begin{center}
        \begin{tabular}{ll}
        \toprule
        直线运动 & 圆周运动 \\
        \midrule
        $x$（位置） & $\theta$（角位置） \\
        $v=\mathrm{d}x/\mathrm{d}t$ & $\omega=\mathrm{d}\theta/\mathrm{d}t$ \\
        $a=\mathrm{d}v/\mathrm{d}t$ & $\alpha=\mathrm{d}\omega/\mathrm{d}t$ \\
        $v=v_{0}+at$ & $\omega=\omega_{0}+\alpha t$ \\
        $x=x_{0}+v_{0}t+\frac{1}{2}at^{2}$ & $\theta=\theta_{0}+\omega_{0}t+\frac{1}{2}\alpha t^{2}$ \\
        $v^{2}-v_{0}^{2}=2a(x-x_{0})$ & $\omega^{2}-\omega_{0}^{2}=2\alpha(\theta-\theta_{0})$ \\
        \bottomrule
        \end{tabular}
        \end{center}

\section{相对运动}
    设参考系 $S'$ 相对参考系 $S$ 作平动（无相对转动），$S'$ 原点对 $S$ 原点的位矢为 $\boldsymbol{r}_{0}$，质点 $P$ 在 $S$、$S'$ 中的位矢分别为 $\boldsymbol{r}$、$\boldsymbol{r}'$，则
    \begin{enumerate}
        \item 位矢变换：$\boldsymbol{r}=\boldsymbol{r}'+\boldsymbol{r}_{0}$；
        \item 位移变换：$\Delta\boldsymbol{r}=\Delta\boldsymbol{r}'+\Delta\boldsymbol{r}_{0}$。
    \end{enumerate}

    \subsection{伽利略速度变换}
    对位矢变换式求时间导数。记质点对 $S$ 的速度为 $\boldsymbol{v}$（绝对速度），对 $S'$ 的速度为 $\boldsymbol{v}'$（相对速度），$S'$ 对 $S$ 的速度为 $\boldsymbol{u}$（牵连速度），则
    \keyformula{\boldsymbol{v}=\boldsymbol{v}'+\boldsymbol{u}}
    这就是\textbf{伽利略速度变换}，其矢量三角形是解题的主要工具。

    \subsection{加速度变换}
    再对时间求导，记 $S'$ 对 $S$ 的加速度为 $\boldsymbol{a}_{0}$：
    \keyformula{\boldsymbol{a}=\boldsymbol{a}'+\boldsymbol{a}_{0}}
    若 $S'$ 相对 $S$ 作匀速直线平动，则 $\boldsymbol{a}_{0}=\boldsymbol{0}$，于是
    \[
    \boldsymbol{a}=\boldsymbol{a}'
    \]
    即所有相对作匀速直线运动的参考系中，测得的质点加速度相同——这正是牛顿定律在其中成立的惯性系的共同特征。

    \subsection{处理相对运动问题的方法}
    \begin{enumerate}
        \item 明确“三个速度”各是谁相对谁，并把下标写成“物体—参考系”的形式（如 $\boldsymbol{v}_{\text{艇岸}}$）；
        \item 画出速度矢量三角形，用正弦定理、余弦定理解三角形；
        \item 或建立坐标系，把矢量方程投影为分量方程求解。
    \end{enumerate}
    \corepoint{伽利略变换（含速度、加速度变换）只在 $v\ll c$ 时成立；高速情形须改用洛伦兹变换（见第5章）。}
